\documentclass[12pt, titlepage]{article}

\usepackage{amsmath, mathtools}

\usepackage[round]{natbib}
\usepackage{amsfonts}
\usepackage{amssymb}
\usepackage{graphicx}
\usepackage{colortbl}
\usepackage{xr}
\usepackage{hyperref}
\usepackage{longtable}
\usepackage{xfrac}
\usepackage{tabularx}
\usepackage{float}
\usepackage{siunitx}
\usepackage{booktabs}
\usepackage{multirow}
\usepackage[section]{placeins}
\usepackage{caption}
\usepackage{fullpage}
\usepackage{listings}

\hypersetup{
bookmarks=true,     % show bookmarks bar?
colorlinks=true,       % false: boxed links; true: colored links
linkcolor=red,          % color of internal links (change box color with linkbordercolor)
citecolor=blue,      % color of links to bibliography
filecolor=magenta,  % color of file links
urlcolor=cyan          % color of external links
}

\usepackage{array}

\externaldocument{../../SRS/SRS}

\input{../../Comments.text}
\input{../../Common.text}

\begin{document}

\title{Module Interface Specification for \progname{}}

\author{\authname}

\date{\today}

\maketitle

\pagenumbering{roman}








% 1 Revision History
\section{Revision History}


\begin{tabularx}{\textwidth}{p{3cm}p{2cm}X}
\toprule {\bf Date} & {\bf Version} & {\bf Notes}\\
\midrule
    Mar 9, 2026 & 1.0 & Initial draft\\
    % Date 2 & 1.1 & Notes\\
\bottomrule
\end{tabularx}


~\newpage








% 2 Symbols, Abbreviations and Acronyms
\section{Symbols, Abbreviations and Acronyms}


See \href{https://zhangwenchili.github.io/CAS741-EelbrainPipeline/SRS/SRS.pdf}{SRS Documentation}.


\newpage


\tableofcontents


\newpage


\pagenumbering{arabic}








% 3 Introduction
\section{Introduction}


The following document details the Module Interface Specifications for \progname{}.


Complementary documents include the System Requirement Specifications
and Module Guide.  The full documentation and implementation can be found at \url{https://github.com/zhangwenchili/CAS741-EelbrainPipeline}.








% 4 Notation
\section{Notation}


The structure of the MIS for modules comes from \citet{HoffmanAndStrooper1995}, with the addition that template modules have been adapted from \cite{GhezziEtAl2003}.  The mathematical notation comes from Chapter 3 of \citet{HoffmanAndStrooper1995}.  For instance, the symbol := is used for a multiple assignment statement and conditional rules follow the form $(c_1 \Rightarrow r_1 | c_2 \Rightarrow r_2 | ... | c_n \Rightarrow r_n )$.


The following table summarizes the primitive data types used by \progname{}. 


\begin{center}
\renewcommand{\arraystretch}{1.2}
\noindent 
\begin{tabular}{l l p{7.5cm}} 
\toprule 
\textbf{Data Type} & \textbf{Notation} & \textbf{Description}\\ 
\midrule
    character & char & a single symbol or digit\\
    integer & $\mathbb{Z}$ & a number without a fractional component in (-$\infty$, $\infty$) \\
    natural number & $\mathbb{N}$ & a number without a fractional component in [1, $\infty$) \\
    real & $\mathbb{R}$ & any number in (-$\infty$, $\infty$)\\
\bottomrule
\end{tabular} 
\end{center}


\noindent The specification of \progname{} \ uses some derived data types: sequences, strings, and tuples. Sequences are lists filled with elements of the same data type. Strings are sequences of characters. Tuples contain a list of values, potentially of different types. In addition, \progname{} uses functions, which are defined by the data types of their inputs and outputs. Local functions are described by giving their type signature followed by their specification.








% 5 Module Decomposition
\section{Module Decomposition}


The following table is taken directly from the Module Guide document for this project.


\begin{table}[h!]
\centering
\begin{tabular}{p{0.3\textwidth} p{0.45\textwidth} p{0.15\textwidth}}
\toprule
\textbf{Level 1} & \textbf{Level 2} & \textbf{Module ID}\\
\midrule

{Hardware-Hiding}
    & Python \\
\midrule

\multirow{7}{0.3\textwidth}{Behaviour-Hiding}
    & TreeModel Module & M2 \\
    & FileTree Module & M3 \\
    & Preprocessing Module & M4 \\
    & Epoch Module & M5 \\
    & Covariance Module & M6 \\
    & MNE-Python & M7 \\
\midrule

\multirow{1}{0.3\textwidth}{Software Decision}
    & Pipeline Module & M1 \\

\bottomrule

\end{tabular}
\caption{Module Hierarchy}
\label{TblMH}
\end{table}

\newpage
~\newpage








% 6 MIS of Module Name
\section{MIS of Pipeline Module}




\subsection{Module}


Pipeline Module




\subsection{Uses}


FileTree Module, Preprocessing Module, Epoch Module, Covariance Module, MNE-Python




\subsection{Syntax}


\subsubsection{Exported Constants}




\subsubsection{Exported Access Programs}


\begin{center}
\begin{tabular}{p{3cm} p{4cm} p{4cm} p{2cm}}
\hline
    \textbf{Name} & \textbf{In} & \textbf{Out} & \textbf{Exceptions} \\
\hline
    \lstinline|load_raw| & - & - & - \\
    \lstinline|load_events| & - & - & - \\
    \lstinline|load_epochs| & - & - & - \\
    \lstinline|load_evoked| & - & - & - \\
    \lstinline|load_evoked_stc| & - & - & - \\
\hline
\end{tabular}
\end{center}




\subsection{Semantics}


\subsubsection{State Variables}


subject, session, task, run, raw, epoch, cov




\subsubsection{Environment Variables}


Eelbrain cache and MRI files in dataset




\subsubsection{Assumptions}


None




\subsubsection{Access Routine Semantics}


\noindent \lstinline|load_raw()|:
\begin{itemize}
    \item output:
    \item exception:
\end{itemize}


\noindent \lstinline|load_events()|:
\begin{itemize}
    \item output:
    \item exception:
\end{itemize}


\noindent \lstinline|load_epochs()|:
\begin{itemize}
    \item output:
    \item exception:
\end{itemize}


\noindent \lstinline|load_evoked()|:
\begin{itemize}
    \item output:
    \item exception:
\end{itemize}


\noindent \lstinline|load_evoked_stc()|:
\begin{itemize}
    \item output:
    \item exception:
\end{itemize}




\subsubsection{Local Functions}


None


\newpage








% 7 MIS of TreeModel Name
\section{MIS of TreeModel Module (M2)}




\subsection{Module}


TreeModel Module




\subsection{Uses}


None




\subsection{Syntax}


\subsubsection{Exported Constants}


None




\subsubsection{Exported Access Programs}


\begin{center}
\begin{tabular}{p{3cm} p{4cm} p{4cm} p{2cm}}
\hline
    \textbf{Name} & \textbf{In} & \textbf{Out} & \textbf{Exceptions} \\
\hline
    \lstinline|_register_field()| & - & - & - \\
    \lstinline|get()| & - & - & - \\
    \lstinline|set()| & - & - & - \\
    \lstinline|iter()| & - & - & - \\
    \lstinline|get_field_values()| & - & - & - \\
\hline
\end{tabular}
\end{center}




\subsection{Semantics}


\subsubsection{State Variables}


Any fileds registered




\subsubsection{Environment Variables}


None




\subsubsection{Assumptions}


None




\subsubsection{Access Routine Semantics}


\noindent \lstinline|_register_field()|:
\begin{itemize}
    \item output:
    \item exception:
\end{itemize}


\noindent \lstinline|get()|:
\begin{itemize}
    \item output:
    \item exception:
\end{itemize}


\noindent \lstinline|set()|:
\begin{itemize}
    \item output:
    \item exception:
\end{itemize}


\noindent \lstinline|iter()|:
\begin{itemize}
    \item output:
    \item exception:
\end{itemize}


\noindent \lstinline|get_field_values()|:
\begin{itemize}
    \item output:
    \item exception:
\end{itemize}




\subsubsection{Local Functions}


None


\newpage








% 8 MIS of FileTree Name
\section{MIS of FileTree Module (M3)}




\subsection{Module}


FileTree Module




\subsection{Uses}


TreeModel Module




\subsection{Syntax}


\subsubsection{Exported Constants}


None




\subsubsection{Exported Access Programs}


\begin{center}
\begin{tabular}{p{3cm} p{4cm} p{4cm} p{2cm}}
\hline
    \textbf{Name} & \textbf{In} & \textbf{Out} & \textbf{Exceptions} \\
\hline
    \lstinline|get()| & - & - & - \\
    \lstinline|glob()| & - & - & - \\
\hline
\end{tabular}
\end{center}




\subsection{Semantics}


\subsubsection{State Variables}


None




\subsubsection{Environment Variables}


None




\subsubsection{Assumptions}


None




\subsubsection{Access Routine Semantics}


\noindent \lstinline|get()|:
\begin{itemize}
    \item output:
    \item exception:
\end{itemize}


\noindent \lstinline|glob()|:
\begin{itemize}
    \item output:
    \item exception:
\end{itemize}




\subsubsection{Local Functions}


None


\newpage








\bibliographystyle {plainnat}
\bibliography {../../../refs/References}


\newpage


\section{Appendix} \label{Appendix}


\newpage{}








\end{document}